\documentclass[11pt]{article}

\usepackage[margin=1in]{geometry}
\usepackage{amsmath,amssymb,amsthm,mathtools}
\usepackage{booktabs,tabularx,array,enumitem,microtype}
\usepackage[hidelinks]{hyperref}
\usepackage[nameinlink,capitalise,noabbrev]{cleveref}
\usepackage{cleveref}
\usepackage{aliascnt}

\newtheorem{theorem}{Theorem}[section]
\newaliascnt{proposition}{theorem}
\newtheorem{proposition}[proposition]{Proposition}
\aliascntresetthe{proposition}
\newaliascnt{lemma}{theorem}
\newtheorem{lemma}[lemma]{Lemma}
\aliascntresetthe{lemma}
\newaliascnt{corollary}{theorem}
\newtheorem{corollary}[corollary]{Corollary}
\aliascntresetthe{corollary}
\newaliascnt{remark}{theorem}
\newtheorem{remark}[remark]{Remark}
\aliascntresetthe{remark}
\theoremstyle{definition}
\newaliascnt{definition}{theorem}
\newtheorem{definition}[definition]{Definition}
\aliascntresetthe{definition}

\crefname{theorem}{theorem}{theorems}
\Crefname{theorem}{Theorem}{Theorems}
\crefname{proposition}{proposition}{propositions}
\Crefname{proposition}{Proposition}{Propositions}
\crefname{lemma}{lemma}{lemmas}
\Crefname{lemma}{Lemma}{Lemmas}
\crefname{corollary}{corollary}{corollaries}
\Crefname{corollary}{Corollary}{Corollaries}
\crefname{remark}{remark}{remarks}
\Crefname{remark}{Remark}{Remarks}

\newcommand{\R}{\mathbb R}
\newcommand{\Pp}{\mathbb P}
\newcommand{\Ee}{\mathbb E}
\newcommand{\cN}{\mathcal N}
\newcommand{\Id}{I}
\newcommand{\tr}{\operatorname{tr}}
\newcommand{\Var}{\operatorname{Var}}
\newcommand{\Cov}{\operatorname{Cov}}
\newcommand{\ip}[2]{\left\langle #1,#2\right\rangle_F}
\newcommand{\norm}[1]{\left\lVert #1\right\rVert_2}
\newcommand{\normF}[1]{\left\lVert #1\right\rVert_F}
\newcommand{\normop}[1]{\left\lVert #1\right\rVert_{\mathrm{op}}}
\newcommand{\normnuc}[1]{\left\lVert #1\right\rVert_*}
\newcommand{\RF}{R_{\mathrm F}}
\newcommand{\Rop}{R_{\mathrm{op}}}
\newcommand{\eps}{\varepsilon}

\title{Sharp Accuracy-Dependent Frobenius--Operator Radius Separation\\
for Gaussian Associative Memory with Zipf Frequencies}
\author{Technical note}
\date{August 2026}

\begin{document}
\maketitle

\begin{abstract}
Let $N=\lceil d^\beta\rceil$ and let the memory frequencies satisfy
$p_i\propto i^{-\alpha}$ for fixed $\alpha\geq0$.  In the
superquadratic regime $\beta>2$, we identify explicit loss intervals on
which the minimum Frobenius and operator radii have matching asymptotic
separation.  In the subquadratic regime $0<\beta<2$, the loss has
infimum zero, and throughout the sharp full-memory interpolation phase
$0<\eps\lesssim Np_N\log N$ we prove the subquadratic separation is $\Theta(\sqrt{\min\{N,d\}})$.  For every fixed $q>0$, all positive
accuracy-dependent conclusions hold simultaneously with probability at
least $1-Cd^{-q}-Ce^{-cd}$.
\end{abstract}

% ============================================================
\section{Model and main theorem}
% ============================================================

We first treat the superquadratic regime.  Until
\Cref{sec:subquadratic}, assume
\[
 \alpha\geq0,\qquad \beta>2,\qquad
 N:=\lceil d^\beta\rceil,\qquad h:=\log N.
\]
The complementary regime $0<\beta<2$ is treated separately in
\Cref{sec:subquadratic}; the critical case $\beta=2$ is not claimed in
this note.  Draw two mutually independent families
\[
 u_1,\ldots,u_N,\ v_1,\ldots,v_N
 \overset{\mathrm{ind}}{\sim}\cN(0,\Id_d/d).
\]
Set
\[
 p_i:=\frac{i^{-\alpha}}{H_{N,\alpha}},
 \qquad H_{N,\alpha}:=\sum_{k=1}^N k^{-\alpha}.
\]
For $W\in\R^{d\times d}$ define
\[
 \ell_i(W):=
 \log\sum_{j=1}^N
 \exp\!\bigl((u_j-u_i)^\top Wv_i\bigr),
 \qquad
 L(W):=\sum_{i=1}^N p_i\ell_i(W).
\]
Thus $L(0)=h$.  For $X\in\{\mathrm F,\mathrm{op}\}$ let
\[
 R_X(\eps):=\inf\{\|W\|_X:L(W)\leq\eps\},
 \qquad \inf\varnothing:=+\infty,
\]
and write
\[
 L_\star:=\inf_W L(W),\qquad
 \Delta:=h-\eps,\qquad
 \mathcal R_X(\Delta):=R_X(h-\Delta).
\]

All limits are as $d\to\infty$ with $(\alpha,\beta)$ fixed.  After a
failure exponent $q>0$ is fixed, every implicit constant may depend on
$(\alpha,\beta,q)$ but not on $d,N,\eps,\Delta$, or an auxiliary head
size $K$.  We write $a\lesssim b$ and $a\asymp b$ in this sense.  An
interval statement
\[
 A_d\lesssim x\lesssim B_d
\]
means that there are fixed constants $C,c>0$ such that the conclusion
holds for every $CA_d\leq x\leq cB_d$.  No relation between $C$ and $c$
is intended; nonemptiness is verified separately whenever needed.

\begin{lemma}[Feasible levels and the universal ratio]
\label{lem:feasible}
If $\eps<L_\star$, then $R_X(\eps)=+\infty$, while $R_X(h)=0$.  If
$L_\star<\eps<h$, or if $\eps=L_\star<h$ and the infimum is attained by
a finite matrix, then both radii are finite and strictly positive and
\begin{equation}
 1\leq \frac{\RF(\eps)}{\Rop(\eps)}\leq\sqrt d.
 \label{eq:universal-ratio}
\end{equation}
\end{lemma}

\begin{proof}
Only positivity needs proof.  Since $L$ is continuous and $L(0)=h>\eps$,
there is $r>0$ such that $\normF{W}<r$ implies $L(W)>\eps$.  Hence
$\RF(\eps)\geq r$ and, because $\normF{W}\leq\sqrt d\normop{W}$,
$\Rop(\eps)\geq r/\sqrt d$.  The norm inequalities
$\normop{W}\leq\normF{W}\leq\sqrt d\normop{W}$ give
\eqref{eq:universal-ratio} after taking infima over the common feasible
set.  The remaining claims follow directly from the definitions and
$L(0)=h$.
\end{proof}

\begin{theorem}[Unified accuracy-dependent separation]
\label{thm:main}
Fix $q>0$.  For all sufficiently large $d$, there is an event of
probability at least
\begin{equation}
 1-Cd^{-q}-Ce^{-cd}
 \label{eq:main-prob}
\end{equation}
on which the following statements hold simultaneously for every loss
level in the displayed interval.
\begin{enumerate}[label=\textup{(\roman*)}]
\item If $0\leq\alpha<1/2$, then
\[
 0<\Delta\lesssim \frac{d^2}{N}
 \quad\Longrightarrow\quad
 \frac{\mathcal R_F(\Delta)}{\mathcal R_{\mathrm{op}}(\Delta)}
 \asymp\sqrt d.
\]
\item If $\alpha=1/2$, then
\[
 0<\Delta\lesssim \frac{d^2h}{N}
 \quad\Longrightarrow\quad
 \frac{\mathcal R_F(\Delta)}{\mathcal R_{\mathrm{op}}(\Delta)}
 \asymp\sqrt d.
\]
\item If $1/2<\alpha<1$, then
\[
 0<\Delta\lesssim N^{2\alpha-2}
 \quad\Longrightarrow\quad
 \frac{\mathcal R_F(\Delta)}{\mathcal R_{\mathrm{op}}(\Delta)}
 \asymp d^{1-\alpha}.
\]
If additionally
\begin{equation}
 \beta<\frac{2-\alpha}{1-\alpha},
 \label{eq:between-condition}
\end{equation}
then
\begin{equation}
 h\left(\frac dN\right)^{1-\alpha}
 \lesssim\Delta\lesssim
 h\left(\frac{d^2}{Nh}\right)^{(1-\alpha)/\alpha}
 \quad\Longrightarrow\quad
 \frac{\mathcal R_F(\Delta)}{\mathcal R_{\mathrm{op}}(\Delta)}
 \asymp\sqrt d.
 \label{eq:between-head-main}
\end{equation}
The interval in \eqref{eq:between-head-main} is nonempty for all large
$d$.
\item If $\alpha=1$, then
\[
 0<\Delta\lesssim h^{-2}
 \quad\Longrightarrow\quad
 \frac{\mathcal R_F(\Delta)}{\mathcal R_{\mathrm{op}}(\Delta)}
 \asymp h.
\]
Moreover, for a fixed $A=A(\beta,q)>0$,
\begin{equation}
 \log\!\left(\frac{Nh^2}{d^2}\right)+A
 \leq\eps\leq
 \log\!\left(\frac Nd\right)-A
 \quad\Longrightarrow\quad
 \frac{\RF(\eps)}{\Rop(\eps)}\asymp\sqrt d.
 \label{eq:alpha-one-main}
\end{equation}
\item If $\alpha>1$, then
\begin{equation}
 \eps\gtrsim\frac{h^\alpha}{d^{2(\alpha-1)}},\qquad \eps<h
 \quad\Longrightarrow\quad
 \frac{\RF(\eps)}{\Rop(\eps)}
 \asymp
 \min\left\{\sqrt d,
 \left(\frac h\eps\right)^{1/(2(\alpha-1))}\right\}.
 \label{eq:alpha-greater-main}
\end{equation}
In particular,
\[
 \frac{h^\alpha}{d^{2(\alpha-1)}}\lesssim\eps
 \lesssim\frac h{d^{\alpha-1}}
 \quad\Longrightarrow\quad
 \frac{\RF(\eps)}{\Rop(\eps)}\asymp\sqrt d.
\]
This last interval has additive width $\asymp h/d^{\alpha-1}$ and
multiplicative endpoint ratio $\asymp(d/h)^{\alpha-1}$.
\end{enumerate}

Finally, suppose $1/2<\alpha<1$ and $\beta(1-\alpha)>2$.  Let
$\mathcal E_{\min}$ be the finite-minimizer event from the companion
note.  On the intersection with the event above, simultaneously for all
feasible $L_\star\leq\eps<h$,
\begin{equation}
 \frac{\RF(\eps)}{\Rop(\eps)}\lesssim d^{1-\alpha}.
 \label{eq:weak-main}
\end{equation}
The probability of this intersection is at least
\begin{equation}
 1-Cd^{-q}-Ce^{-cd}
 -2^{1-N}\sum_{r=0}^{d^2-1}\binom{N-1}{r}.
 \label{eq:weak-main-prob}
\end{equation}
\end{theorem}

\paragraph{Proof structure and probability bookkeeping.}
The proof has four modules: deterministic convex geometry; a base
Gaussian event controlling sample covariances and the initial gradient;
a diffuse event for $\alpha\leq1/2$ or a head event for $\alpha>1/2$;
and, only for \eqref{eq:weak-main}, the companion minimizer event.
Every Gaussian lemma below is simultaneous over all indices used later.
The diffuse and head arguments use union bounds over at most $N^3$
choices, hence contribute $Cd^{-q}$; Gaussian norm and spectral events
contribute $Ce^{-cd}$.  Once these events occur, all statements over a
continuum of loss levels are deterministic consequences of convex
scaling or interval coverage, so no union bound over $\eps$ is used.

% ============================================================
\section{Deterministic convex geometry}
% ============================================================

Set
\[
 \bar u:=\frac1N\sum_{j=1}^N u_j,
 \qquad x_j:=u_j-\bar u,
 \qquad G:=\sum_{i=1}^N p_i x_i v_i^\top.
\]
Then $\sum_jx_j=0$.

\begin{lemma}[Centered representation]
\label{lem:centered}
For every $W$,
\begin{equation}
 L(W)=h+\Psi(W)-\ip{G}{W},
 \qquad
 \Psi(W):=\sum_{i=1}^N p_i
 \log\left(\frac1N\sum_{j=1}^N e^{x_j^\top Wv_i}\right).
 \label{eq:centered}
\end{equation}
Moreover,
\[
 \Psi(0)=0,\qquad \nabla\Psi(0)=0,\qquad \Psi(W)\geq0,
\]
and, for every direction $A$,
\begin{equation}
 D^2\Psi(W)[A,A]
 =\sum_{i=1}^N p_i
 \Var_{j\sim q_i(W)}(x_j^\top Av_i),
 \label{eq:hessian-variance}
\end{equation}
where
\[
 q_{ji}(W):=\frac{e^{x_j^\top Wv_i}}
 {\sum_{k=1}^N e^{x_k^\top Wv_i}}.
\]
\end{lemma}

\begin{proof}
Since $u_j-u_i=x_j-x_i$,
\[
 \ell_i(W)=h+
 \log\left(\frac1N\sum_j e^{x_j^\top Wv_i}\right)
 -x_i^\top Wv_i,
\]
which gives \eqref{eq:centered}.  Jensen's inequality and
$N^{-1}\sum_jx_j=0$ give $\Psi(W)\geq0$, and differentiation at zero
gives $\nabla\Psi(0)=0$.  A second differentiation of each
log-partition function gives the variance formula
\eqref{eq:hessian-variance}.
\end{proof}

\begin{proposition}[Tangent bounds, local response, and scaling]
\label{prop:convex-geometry}
Let $\Delta>0$.
\begin{enumerate}[label=\textup{(\roman*)}]
\item Every $W$ with $L(W)\leq h-\Delta$ satisfies
\[
 \ip{G}{W}\geq\Delta.
\]
If $G=0$, the positive-gain sublevel set is therefore empty.  If
$G\neq0$, then
\begin{equation}
 \mathcal R_F(\Delta)\geq\frac{\Delta}{\normF{G}},
 \qquad
 \mathcal R_{\mathrm{op}}(\Delta)\geq
 \frac{\Delta}{\normnuc{G}},
 \label{eq:tangent}
\end{equation}
with the usual convention if the sublevel set is empty.
\item Assume
\begin{equation}
 \max_j\norm{x_j}\leq3,
 \qquad
 \sum_i p_i\norm{v_i}^2\leq4.
 \label{eq:deterministic-regularity}
\end{equation}
Then
\begin{equation}
 \Psi(A)\leq18\normop{A}^2
 \label{eq:global-psi}
\end{equation}
for every $A$.  Hence, whenever
$0<\Delta\leq\normF{G}^2/72$,
\begin{equation}
 \mathcal R_F(\Delta)\asymp\frac{\Delta}{\normF{G}},
 \qquad
 \mathcal R_{\mathrm{op}}(\Delta)
 \asymp\frac{\Delta}{\normnuc{G}}.
 \label{eq:local-radii}
\end{equation}
\item If $L(W_0)\leq h-\Delta_0$, then
$(\Delta/\Delta_0)W_0$ is feasible at gain $\Delta$ for every
$0\leq\Delta\leq\Delta_0$.
\item If a finite minimizer $W_\star$ exists and
$\Delta_\star:=h-L(W_\star)>0$, then, for every
$0<\Delta\leq\Delta_\star$,
\begin{equation}
 \frac{\mathcal R_F(\Delta)}
 {\mathcal R_{\mathrm{op}}(\Delta)}
 \leq
 \frac{\normF{W_\star}\normnuc{G}}{\Delta_\star}.
 \label{eq:all-accuracy}
\end{equation}
\end{enumerate}
\end{proposition}

\begin{proof}
Part (i) follows from \eqref{eq:centered} and $\Psi\geq0$, followed by
Frobenius and nuclear--operator duality.

For part (ii), \eqref{eq:hessian-variance} and
\eqref{eq:deterministic-regularity} give, for every $W,A$,
\[
 D^2\Psi(W)[A,A]
 \leq\sum_i p_i\max_j|x_j^\top Av_i|^2
 \leq36\normop{A}^2.
\]
Taylor's formula at zero yields \eqref{eq:global-psi}.  For the
Frobenius upper bound, take $A=G/\normF{G}$ and
$t=2\Delta/\normF{G}$; then
\[
 L(tA)\leq h-2\Delta+72\frac{\Delta^2}{\normF{G}^2}
 \leq h-\Delta.
\]
For the operator upper bound, write $G=U\Sigma V^\top$ and use
$Q=UV^\top$, for which $\normop{Q}=1$ and
$\ip{G}{Q}=\normnuc{G}$.  The same calculation applies because
$\normnuc{G}\geq\normF{G}$.  Together with \eqref{eq:tangent}, this
proves \eqref{eq:local-radii}.

Part (iii) is convexity:
\[
 L\left(\frac\Delta{\Delta_0}W_0\right)
 \leq\left(1-\frac\Delta{\Delta_0}\right)L(0)
 +\frac\Delta{\Delta_0}L(W_0)
 \leq h-\Delta.
\]
For part (iv), scale $W_\star$ as in part (iii) to upper-bound
$\mathcal R_F(\Delta)$ and combine with the operator tangent bound in
\eqref{eq:tangent}.
\end{proof}

% ============================================================
\section{Zipf sums and Gaussian estimates}
% ============================================================

We repeatedly use the following standard estimates
\cite{vershynin,janson}.  If $g_r$ are independent standard Gaussians
and $b_r\geq0$ are deterministic, then
\begin{equation}
 \Pp\!\left(
 \left|\sum_r b_r(g_r^2-1)\right|
 >2\sqrt{t\sum_r b_r^2}+2t\max_r b_r
 \right)\leq2e^{-t}.
 \label{eq:weighted-chi-square}
\end{equation}
If $Y$ is an $m\times n$ standard Gaussian matrix with $n\geq m$, then
\begin{equation}
 \Pp\!\left(
 s_{\max}(Y)>\sqrt n+\sqrt m+t
 \ \text{or}\ 
 s_{\min}(Y)<\sqrt n-\sqrt m-t
 \right)\leq2e^{-t^2/2}.
 \label{eq:gaussian-singular-values}
\end{equation}
If $f$ is $L$-Lipschitz and $g\sim\cN(0,\Id_d/d)$, then
\begin{equation}
 \Pp(|f(g)-\Ee f(g)|>t)\leq
 2\exp\!\left(-\frac{dt^2}{2L^2}\right).
 \label{eq:gaussian-lipschitz}
\end{equation}
If $X_k$ are independent and centered, then
\begin{equation}
 \Pp\!\left(
 \left|\sum_kX_k\right|>
 C\left[
 \sqrt{t\sum_k\|X_k\|_{\psi_1}^2}
 +t\max_k\|X_k\|_{\psi_1}
 \right]
 \right)\leq2e^{-t}.
 \label{eq:subexponential-bernstein}
\end{equation}
Finally, a Gaussian polynomial of degree at most four satisfies
$\|Z\|_{L^r}\leq(r-1)^2\|Z\|_{L^2}$ for $r\geq2$.

\begin{lemma}[Zipf estimates]
\label{lem:zipf}
Let
\[
 s_2:=\sum_{i=1}^N p_i^2,
 \qquad
 P_{>K}:=\sum_{i>K}p_i,
 \qquad
 s_{2,>K}:=\sum_{i>K}p_i^2.
\]
Then
\begin{equation}
 s_2\asymp
 \begin{cases}
 N^{-1},&0\leq\alpha<1/2,\\
 h/N,&\alpha=1/2,\\
 N^{2\alpha-2},&1/2<\alpha<1,\\
 h^{-2},&\alpha=1,\\
 1,&\alpha>1.
 \end{cases}
 \label{eq:s2-orders}
\end{equation}
Moreover, uniformly for $1\leq K\leq N/2$,
\begin{align}
 0\leq\alpha<1:
 &\quad p_K\asymp N^{\alpha-1}K^{-\alpha},
 \quad \sum_{i\leq K}p_i\asymp Kp_K
 \asymp(K/N)^{1-\alpha};
 \label{eq:zipf-below-one}\\
 1/2<\alpha<1:
 &\quad s_{2,>K}\asymp N^{2\alpha-2}K^{1-2\alpha};
 \label{eq:zipf-tail-between}\\
 \alpha=1:
 &\quad p_K\asymp(Kh)^{-1},
 \quad P_{>K}h=\log(N/K)+O(1),
 \quad s_{2,>K}\asymp(Kh^2)^{-1};
 \label{eq:zipf-one}\\
 \alpha>1:
 &\quad p_K\asymp K^{-\alpha},
 \quad P_{>K}\asymp K^{1-\alpha},
 \quad s_{2,>K}\asymp K^{1-2\alpha}.
 \label{eq:zipf-above-one}
\end{align}
At $\alpha=1/2$,
$s_{2,>K}\asymp N^{-1}\log(N/K)$.
\end{lemma}

\begin{proof}
All statements follow from the integral comparisons
\[
 \sum_{i=1}^M i^{-r}\asymp
 \begin{cases}
 M^{1-r},&r<1,\\
 \log M,&r=1,\\
 1,&r>1,
 \end{cases}
\]
and, for $r>1$ and $K\leq M/2$,
$\sum_{i=K+1}^M i^{-r}\asymp K^{1-r}$, after dividing by
$H_{N,\alpha}$ or $H_{N,\alpha}^2$ as appropriate.  For
$\alpha=1$, use
$H_{N,1}=h+O(1)$ and
$H_{N,1}-H_{K,1}=\log(N/K)+O(K^{-1})$.
\end{proof}

\begin{lemma}[Gaussian bilinear sums]
\label{lem:bilinear}
Let $a_1,\ldots,a_m$ be deterministic, let
$u_i,v_i\overset{\mathrm{ind}}\sim\cN(0,\Id_d/d)$, and put
\[
 Z:=\sum_{i=1}^m a_i u_i v_i^\top,
 \quad A_2:=\left(\sum_i a_i^2\right)^{1/2},
 \quad A_\infty:=\max_i|a_i|.
\]
Then, with probability at least $1-Ce^{-cd}$,
\begin{equation}
 \normF{Z}\asymp A_2,
 \qquad
 \normop{Z}\lesssim A_\infty+\frac{A_2}{\sqrt d}.
 \label{eq:bilinear-matrix}
\end{equation}
The upper bounds hold simultaneously for any polynomial number of
deterministic coefficient vectors.

If $Q$ is deterministic and
$X=\sum_i a_i u_i^\top Qv_i$, then for every $t\geq1$, with
probability at least $1-2e^{-t}$,
\begin{equation}
 |X|\lesssim
 \frac{\normF{Q}A_2\sqrt t}{d}
 +\frac{\normop{Q}A_\infty t}{d}.
 \label{eq:scalar-bilinear}
\end{equation}
The same statement holds conditionally if $Q$ is random but independent
of the Gaussian pairs in the sum.
\end{lemma}

\begin{proof}
Write $U=[u_1,\ldots,u_m]$, $V=[v_1,\ldots,v_m]$, and
$D=\operatorname{diag}(a_1,\ldots,a_m)$.  Formula
\eqref{eq:weighted-chi-square}, first for $UD$ and then conditionally on
$U$ for the columns of $UDV^\top$, gives
$\normF{UD}\asymp A_2$ and
$\normF{UDV^\top}\asymp\normF{UD}$, each with failure $Ce^{-cd}$.
A $1/4$-net and \eqref{eq:weighted-chi-square} give
\[
 \normop{UD}\lesssim A_\infty+A_2/\sqrt d.
\]
Conditionally on $U$,
$UDV^\top$ has the same law as
$(UD^2U^\top)^{1/2}Y/\sqrt d$, where $Y$ is a standard $d\times d$
Gaussian matrix.  Since $\normop{Y}\lesssim\sqrt d$ with failure
$Ce^{-cd}$, this proves \eqref{eq:bilinear-matrix}.

For the scalar statement, singular-value decomposition and rotational
invariance give, for
$|\lambda|A_\infty\normop{Q}/d\leq1/2$,
\[
 \log\Ee e^{\lambda X}
 \leq \lambda^2 A_2^2\normF{Q}^2/d^2.
\]
Chernoff's method, applied also to $-X$, yields
\eqref{eq:scalar-bilinear}.
\end{proof}

\begin{proposition}[Base Gaussian event and initial gradient]
\label{prop:base-event}
Fix $q>0$.  With probability at least
$1-Cd^{-q}-Ce^{-cd}$, the following estimates hold simultaneously:
\begin{align}
 &\max_i\bigl|\norm{u_i}^2-1\bigr|
 +\max_i\bigl|\norm{v_i}^2-1\bigr|\leq\frac1{10},
 \label{eq:base-norms}\\
 &\max_{i\neq j}
 \bigl(|u_i^\top u_j|+|v_i^\top v_j|\bigr)
 \lesssim\sqrt{h/d},
 \label{eq:base-inner}\\
 &\norm{\bar u}\lesssim N^{-1/2},
 \qquad
 \frac c d\Id_d\preceq
 \Sigma_x:=\frac1N\sum_{j=1}^N x_jx_j^\top
 \preceq\frac C d\Id_d,
 \label{eq:base-Sigma}\\
 &\normF{G}\asymp\sqrt{s_2}.
 \label{eq:base-GF}
\end{align}
If $\alpha\leq1/2$, then
\begin{equation}
 M_v:=\sum_i p_i v_iv_i^\top\preceq C\Id_d/d.
 \label{eq:base-Mv-upper}
\end{equation}
If $1/2<\alpha<1$ and $\beta(1-\alpha)>2$, then
\begin{equation}
 c\Id_d/d\preceq M_v\preceq C\Id_d/d.
 \label{eq:base-Mv-two-sided}
\end{equation}
Finally, for $\alpha>1/2$,
\begin{equation}
 \frac{\normnuc{G}}{\normF{G}}
 \asymp
 \begin{cases}
 d^{1-\alpha},&1/2<\alpha<1,\\
 h,&\alpha=1,\\
 1,&\alpha>1.
 \end{cases}
 \label{eq:base-nuclear}
\end{equation}
\end{proposition}

\begin{proof}
The Gaussian norm and inner-product tails, followed by union bounds over
$O(N^2)$ indices, give \eqref{eq:base-norms}--\eqref{eq:base-inner}.
The singular-value bounds for the $d\times N$ matrix
$[u_1,\ldots,u_N]$ give two-sided bounds of order $d^{-1}$ for its
sample covariance.  Since
\[
 \Sigma_x=\frac1N\sum_j u_ju_j^\top-\bar u\bar u^\top,
 \qquad \norm{\bar u}^2=O_{\Pp}(N^{-1})=o(d^{-1}),
\]
this proves \eqref{eq:base-Sigma}.

For fixed $z\in\mathbb S^{d-1}$,
\[
 d z^\top M_vz-1=\sum_i p_i(g_i^2-1).
\]
By \eqref{eq:weighted-chi-square}, a fixed-direction deviation by a
constant has exponent
$\gtrsim\min\{s_2^{-1},p_1^{-1}\}$.  For $\alpha\leq1/2$ this exponent
is $\gg d$ because $\beta>2$; in the weak regime it is $\gg d$ because
$p_1^{-1}\asymp N^{1-\alpha}$ and $\beta(1-\alpha)>2$.  A $1/4$-net
therefore gives \eqref{eq:base-Mv-upper} and
\eqref{eq:base-Mv-two-sided}.

Write
\[
 \widetilde G:=\sum_i p_i u_iv_i^\top,
 \qquad z_p:=\sum_i p_i v_i,
 \qquad G=\widetilde G-\bar u z_p^\top.
\]
By \Cref{lem:bilinear},
$\normF{\widetilde G}\asymp\sqrt{s_2}$.  Also
$\norm{\bar u}\lesssim N^{-1/2}$ and
$\norm{z_p}\lesssim\sqrt{s_2}$, so the centering term is
$o(\sqrt{s_2})$.  This proves \eqref{eq:base-GF}.

It remains to prove \eqref{eq:base-nuclear}.  Let
$m=\lfloor d/4\rfloor$, set
$U_m=[u_1,\ldots,u_m]$, $V_m=[v_1,\ldots,v_m]$, and define
\[
 Q_0:=U_m(U_m^\top U_m)^{-1}(V_m^\top V_m)^{-1}V_m^\top.
\]
With failure $Ce^{-cd}$, the singular values of $U_m,V_m$ are bounded
above and below by constants.  Hence, after multiplying $Q_0$ by a fixed
constant, we obtain $\normop{Q}\leq1$ and
$u_i^\top Qv_i\geq c$ for all $i\leq m$.  Conditional on the head data, apply the scalar part of
\Cref{lem:bilinear} with $t=(q+2)h$ to obtain
\[
 \left|\sum_{i>m}p_i u_i^\top Qv_i\right|
 \lesssim
 \sqrt{\frac{s_{2,>m}h}{d}}+\frac{p_{m+1}h}{d}
\]
with failure at most $d^{-q}$.  If $1/2<\alpha<1$, the head signal and
the two displayed errors have respective orders
\[
 N^{\alpha-1}d^{1-\alpha},\qquad
 N^{\alpha-1}d^{-\alpha}\sqrt h,
 \qquad N^{\alpha-1}d^{-\alpha-1}h,
\]
so both errors are $o(1)$ times the signal.  At $\alpha=1$, the signal
is $\asymp1$, while the errors are
$O((d\sqrt h)^{-1})$ and $O(d^{-2})$.  The centering correction
$O(\sqrt{s_2/N})$ is also negligible in both cases.  Therefore
\begin{equation}
 \normnuc{G}\gtrsim
 \begin{cases}
 N^{\alpha-1}d^{1-\alpha},&1/2<\alpha<1,\\
 1,&\alpha=1.
 \end{cases}
 \label{eq:nuclear-lower-head}
\end{equation}

For the matching upper bound, split the uncentered gradient at $d$:
\[
 \normnuc{\widetilde G}
 \leq \sum_{i\leq d}p_i\norm{u_i}\norm{v_i}
 +\sqrt d\normF{\sum_{i>d}p_i u_iv_i^\top}.
\]
Using \Cref{lem:bilinear,lem:zipf} and adding the negligible centering
term gives
\[
 \normnuc{G}\lesssim
 \begin{cases}
 N^{\alpha-1}d^{1-\alpha},&1/2<\alpha<1,\\
 1,&\alpha\geq1.
 \end{cases}
\]
Together with \eqref{eq:nuclear-lower-head} and
\eqref{eq:base-GF} this proves the first two cases of
\eqref{eq:base-nuclear}.  If $\alpha>1$, the nuclear norm dominates the
Frobenius norm, while the same split gives $\normnuc{G}\lesssim1$;
thus the ratio is $\asymp1$.
\end{proof}

\begin{corollary}[Local phase for $\alpha>1/2$]
\label{cor:local-phase}
On the event of \Cref{prop:base-event}, simultaneously for
$0<\Delta\lesssim s_2$,
\[
 \frac{\mathcal R_F(\Delta)}
 {\mathcal R_{\mathrm{op}}(\Delta)}
 \asymp
 \begin{cases}
 d^{1-\alpha},&1/2<\alpha<1,\\
 h,&\alpha=1,\\
 1,&\alpha>1.
 \end{cases}
\]
\end{corollary}

\begin{proof}
The base norm event implies \eqref{eq:deterministic-regularity}, while
\eqref{eq:base-GF} gives $\normF{G}^2\asymp s_2$.  Apply
\eqref{eq:local-radii} and then \eqref{eq:base-nuclear}.
\end{proof}

% ============================================================
\section{Diffuse phase for \texorpdfstring{$\alpha\leq1/2$}{alpha <= 1/2}}
% ============================================================

Assume in this section that $0\leq\alpha\leq1/2$.  For all large $d$,
$d^2<N/2$.  Define
\begin{equation}
 T:=\{d^2+1,\ldots,N\},
 \qquad s_T:=\sum_{i\in T}p_i^2,
 \qquad S:=d^2\sum_{i\in T}p_i x_iv_i^\top.
 \label{eq:diffuse-def}
\end{equation}

\begin{proposition}[Diffuse comparator]
\label{prop:diffuse-comparator}
Fix $q>0$.  With probability at least
$1-Cd^{-q}-Ce^{-cd}$,
\begin{align}
 &s_T\asymp s_2\asymp
 \begin{cases}
 N^{-1},&\alpha<1/2,\\
 h/N,&\alpha=1/2,
 \end{cases}
 \label{eq:diffuse-mass}\\
 &\normF{S}\lesssim d^2\sqrt{s_T},
 \qquad
 \normop{S}\lesssim d^{3/2}\sqrt{s_T},
 \label{eq:diffuse-norms}\\
 &\ip{G}{S}\gtrsim d^2s_T,
 \label{eq:diffuse-alignment}\\
 &\max_{i,j}|x_j^\top Sv_i|
 \lesssim
 d\sqrt{s_Th}+d^2p_{d^2+1}
 +d^2\sqrt{s_T/N}=o(1).
 \label{eq:diffuse-logits}
\end{align}
On the same event,
\begin{equation}
 \frac{\normnuc{G}}{\normF{G}}\asymp\sqrt d.
 \label{eq:diffuse-nuclear}
\end{equation}
\end{proposition}

\begin{proof}
For $\alpha<1/2$, integral comparison gives
\[
 \sum_{i>d^2}i^{-2\alpha}
 \asymp N^{1-2\alpha},
\]
because $d^2/N\to0$.  At $\alpha=1/2$ the same tail equals
$\log(N/d^2)+O(1)\asymp h$.  Division by
$H_{N,\alpha}^2$ proves \eqref{eq:diffuse-mass}.

For the uncentered matrix
\[
 \widetilde S:=d^2\sum_{i\in T}p_i u_iv_i^\top,
\]
the coefficient norms are
\[
 A_2=d^2\sqrt{s_T},
 \qquad A_\infty=d^2p_{d^2+1}.
\]
A direct use of \Cref{lem:zipf} gives
\begin{equation}
 A_\infty\lesssim A_2/\sqrt d,
 \qquad A_\infty=o(1).
 \label{eq:diffuse-coefficients}
\end{equation}
For example, when $\alpha<1/2$,
$dA_\infty^2/A_2^2\lesssim d^{-1}(d^2/N)^{1-2\alpha}$, and at
$\alpha=1/2$ this ratio is $O((dh)^{-1})$.
Thus \Cref{lem:bilinear} gives the bounds in
\eqref{eq:diffuse-norms} for $\widetilde S$.  Since
\[
 S=\widetilde S-d^2\bar u
 \left(\sum_{i\in T}p_iv_i\right)^\top,
\]
the centering term has both Frobenius and operator norm at most
$d^2O(N^{-1/2})O(\sqrt{s_T})$, which is negligible relative to the two
scales in \eqref{eq:diffuse-norms}.

We next prove alignment.  Put
$\widetilde G=\sum_i p_i u_iv_i^\top$.  Expanding gives
\[
 \ip{\widetilde G}{\widetilde S}=D+Z,
\]
where
\[
 D=d^2\sum_{i\in T}p_i^2\norm{u_i}^2\norm{v_i}^2
\]
and
\[
 Z=d^2\sum_{k\in T}\sum_{i\neq k}
 p_ip_k(u_i^\top u_k)(v_i^\top v_k).
\]
On \eqref{eq:base-norms}, $D\geq c d^2s_T$.  In the expansion of
$\Ee Z^2$, a Gaussian expectation vanishes unless the second ordered
pair equals $(i,k)$ or $(k,i)$; in either nonzero case the product of
the $u$- and $v$-expectations is $d^{-2}$.  Therefore
\begin{align*}
 \Ee Z^2
 &\leq\frac1{d^2}
 \left(s_2\,d^4s_T+(d^2s_T)^2\right),
\end{align*}
so, using $s_T\asymp s_2$,
\[
 \frac{\|Z\|_{L^2}}{d^2s_T}\lesssim d^{-1}.
\]
The random variable $Z$ belongs to Gaussian chaos of degree four.
Choose an even $r>q+2$ and use hypercontractivity and Markov's
inequality to obtain
\[
 \Pp\bigl(|Z|>d^2s_T/8\bigr)\leq d^{-q}
\]
for all large $d$.  Restoring the two centering terms changes the inner
product by at most
\[
 O\!\left(d^2\sqrt{s_2s_T/N}\right)=o(d^2s_T).
\]
This proves \eqref{eq:diffuse-alignment}.

It remains to prove the uniform logit bound.  Fix $(i,j)$ and condition
only on the anchor vectors $(u_j,v_i)$.  For
$k\in T\setminus\{i,j\}$, set
\[
 X_k:=d^2p_k(u_j^\top u_k)(v_k^\top v_i).
\]
The $X_k$ are conditionally independent and centered.  On the anchor
norm event,
\[
 \sum_k\|X_k\|_{\psi_1\mid u_j,v_i}^2\lesssim d^2s_T,
 \qquad
 \max_k\|X_k\|_{\psi_1\mid u_j,v_i}
 \lesssim dp_{d^2+1}.
\]  Bernstein's
inequality with $t=(q+3)h$ therefore gives
\[
 \left|\sum_{k\in T\setminus\{i,j\}}
 d^2p_k(u_j^\top u_k)(v_k^\top v_i)\right|
 \lesssim d\sqrt{s_Th}+d p_{d^2+1}h.
\]
By \eqref{eq:diffuse-coefficients} and $h/d=o(1)$, the second term is
bounded by a constant multiple of $d^2p_{d^2+1}$.  The at most two
removed anchor terms obey the same bound on
\eqref{eq:base-norms}--\eqref{eq:base-inner}.  A union bound over
$N^2$ pairs leaves failure at most $d^{-q}$.  Finally, replacing both
$u_j$ and $\widetilde S$ by their centered versions adds at most
\[
 O\!\left(d^2\sqrt{s_T/N}\right).
\]
This proves \eqref{eq:diffuse-logits}.  Its right-hand side tends to zero:
for $\alpha<1/2$ it is bounded by
\[
 C\left(d\sqrt{h/N}+(d^2/N)^{1-\alpha}+d^2/N\right),
\]
and at $\alpha=1/2$ by
\[
 C\left(dh/\sqrt N+d/\sqrt N+d^2\sqrt h/N\right).
\]
All terms vanish because $\beta>2$.

Nuclear--operator duality, \eqref{eq:diffuse-alignment}, and
\eqref{eq:diffuse-norms} give
\[
 \normnuc{G}\geq
 \frac{\ip{G}{S}}{\normop{S}}
 \gtrsim\sqrt d\sqrt{s_T}.
\]
Since $\normnuc{G}\leq\sqrt d\normF{G}$ and
$\normF{G}\asymp\sqrt{s_2}\asymp\sqrt{s_T}$, this proves
\eqref{eq:diffuse-nuclear}.
\end{proof}

\begin{theorem}[Diffuse accuracy interval]
\label{thm:diffuse}
Under $0\leq\alpha\leq1/2$, on the event of
\Cref{prop:base-event,prop:diffuse-comparator}, simultaneously for
\begin{equation}
 0<\Delta\lesssim d^2s_T,
 \label{eq:diffuse-gain}
\end{equation}
one has
\begin{equation}
 \mathcal R_{\mathrm{op}}(\Delta)
 \lesssim\frac{\Delta}{\sqrt d\sqrt{s_T}},
 \qquad
 \mathcal R_F(\Delta)\gtrsim\frac{\Delta}{\sqrt{s_2}},
 \label{eq:diffuse-radii}
\end{equation}
and hence
\[
 \frac{\mathcal R_F(\Delta)}
 {\mathcal R_{\mathrm{op}}(\Delta)}\asymp\sqrt d.
\]
\end{theorem}

\begin{proof}
By \eqref{eq:diffuse-logits}, the whole segment $tS$, $0\leq t\leq1$,
lies in a fixed small-logit ball for large $d$.  There every softmax
weight is within a fixed factor of $1/N$.  Using
\eqref{eq:hessian-variance} and the pairwise-difference representation
of covariance,
\[
 D^2\Psi(tS)[S,S]
 \lesssim \tr(S^\top\Sigma_x S M_v)
 \lesssim d^{-2}\normF{S}^2
 \lesssim d^2s_T.
\]
Together with \eqref{eq:diffuse-alignment}, Taylor's formula gives
\[
 L(tS)\leq h-c t d^2s_T+C t^2d^2s_T.
\]
A sufficiently small fixed $t_0>0$ therefore achieves gain
$\gtrsim d^2s_T$.  Convex scaling and \eqref{eq:diffuse-norms} give the
operator upper bound in \eqref{eq:diffuse-radii}.  The Frobenius lower
bound is \eqref{eq:tangent} and \eqref{eq:base-GF}.  Finally use
$s_T\asymp s_2$ and \eqref{eq:universal-ratio}.
\end{proof}

% ============================================================
\section{Head phase for \texorpdfstring{$\alpha>1/2$}{alpha > 1/2}}
% ============================================================

Assume throughout this section that $\alpha>1/2$.  For $K<N$ set
\[
 P_K:=\sum_{i>K}p_i,
 \qquad s_K:=\sum_{i>K}p_i^2,
 \qquad S_K:=\sum_{i=1}^K u_iv_i^\top,
\]
and define
\begin{equation}
 B_K:=\sum_{i=1}^N
 \min\left\{1,\frac{p_i}{p_K}\right\}x_iv_i^\top.
 \label{eq:BK-def}
\end{equation}
Fix once and for all $c_0:=1/8$.

\begin{lemma}[Pooled Jensen lower bound]
\label{lem:pooled}
For every $W$ and every $i$,
\begin{equation}
 \ip{W}{x_iv_i^\top}\geq h-\ell_i(W).
 \label{eq:single-pooled}
\end{equation}
Consequently, if
\begin{equation}
 L(W)\leq P_Kh+c_0Kp_Kh,
 \label{eq:pooled-threshold}
\end{equation}
then
\begin{equation}
 \ip{W}{B_K}\geq(1-c_0)Kh.
 \label{eq:BK-alignment}
\end{equation}
Moreover, with probability at least $1-Ce^{-cd}$, simultaneously for
all $K\leq N/2$,
\begin{equation}
 \normF{B_K}\lesssim\sqrt K.
 \label{eq:BK-norm}
\end{equation}
Thus every $W$ satisfying \eqref{eq:pooled-threshold} obeys
\begin{equation}
 \normF{W}\gtrsim h\sqrt K.
 \label{eq:head-F-lower}
\end{equation}
\end{lemma}

\begin{proof}
Jensen's inequality applied to all $N$ distractors gives
\[
 \ell_i(W)\geq h+\frac1N\sum_j(u_j-u_i)^\top Wv_i
 =h-x_i^\top Wv_i,
\]
which is \eqref{eq:single-pooled}.  Since
\[
 p_K\min\left\{1,\frac{p_i}{p_K}\right\}
 =\begin{cases}p_K,&i\leq K,\\ p_i,&i>K,\end{cases}
\]
multiplying \eqref{eq:single-pooled} by these coefficients and summing
gives
\[
 p_K\ip{W}{B_K}
 \geq(Kp_K+P_K)h-L(W).
\]
Use \eqref{eq:pooled-threshold} and divide by $p_K$.

For the norm bound, the squared coefficient norm in \eqref{eq:BK-def}
is
\[
 K+\sum_{i>K}(p_i/p_K)^2
 =K+K^{2\alpha}\sum_{i>K}i^{-2\alpha}\asymp K,
\]
uniformly for $K\leq N/2$.  Apply \Cref{lem:bilinear} and note that
the centering term has Frobenius norm $O(\sqrt{K/N})$.  A union bound
over $K$ preserves failure $Ce^{-cd}$.  Finally combine
\eqref{eq:BK-alignment} and \eqref{eq:BK-norm} by Cauchy--Schwarz.
\end{proof}

Define the deterministic error scale
\begin{equation}
 \mathsf e_K:=
 P_K\left(\frac{h^2}{d}+\frac{Kh^2}{d^2}\right)
 +h\left(1+\sqrt{K/d}\right)\sqrt{\frac{s_Kh}{d}}
 +N^{-4\alpha-4}.
 \label{eq:head-error}
\end{equation}

\begin{proposition}[Uniform head comparator]
\label{prop:head-comparator}
Fix $q>0$.  There exist fixed constants
\[
 \kappa=\kappa(\beta,q)>0,
 \qquad C_{\mathrm{mem}}=C_{\mathrm{mem}}(\alpha)>0,
 \qquad C_{\mathrm{cmp}}=C_{\mathrm{cmp}}(\alpha,\beta,q)>0,
\]
such that, with probability at least $1-Cd^{-q}-Ce^{-cd}$, the
following statements hold simultaneously for every integer $K$ with
\begin{equation}
 Kh\leq\kappa d^2.
 \label{eq:head-size}
\end{equation}
One has
\begin{align}
 &\normF{S_K}\lesssim\sqrt K,
 \qquad
 \normop{S_K}\lesssim1+\sqrt{K/d},
 \label{eq:SK-norms}\\
 &(u_i-u_j)^\top S_Kv_i\geq\frac14
 \quad(i\leq K,\ j\neq i).
 \label{eq:head-margin}
\end{align}
For
\begin{equation}
 W_K:=C_{\mathrm{mem}}hS_K,
 \label{eq:WK-def}
\end{equation}
one has
\begin{align}
 &\normF{W_K}\lesssim h\sqrt K,
 \qquad
 \normop{W_K}\lesssim h(1+\sqrt{K/d}),
 \label{eq:WK-norms}\\
 &L(W_K)\leq P_Kh+C_{\mathrm{cmp}}\mathsf e_K.
 \label{eq:WK-loss}
\end{align}
If in addition
\begin{equation}
 C_{\mathrm{cmp}}\mathsf e_K
 \leq\frac{c_0}{4}Kp_Kh,
 \label{eq:head-admissible}
\end{equation}
then every
\begin{equation}
 P_Kh+C_{\mathrm{cmp}}\mathsf e_K
 \leq\eps\leq P_Kh+c_0Kp_Kh
 \label{eq:one-head-window}
\end{equation}
satisfies
\begin{equation}
 \Rop(\eps)\lesssim h(1+\sqrt{K/d}),
 \qquad
 \RF(\eps)\gtrsim h\sqrt K.
 \label{eq:one-head-radii}
\end{equation}
\end{proposition}

\begin{proof}
The norm estimates in \eqref{eq:SK-norms} are
\eqref{eq:bilinear-matrix} with $K$ unit coefficients, followed by a
union bound over $K\leq N$.

For the margin, fix $(K,i,j)$ with $i\leq K$ and $j\neq i$.  Decompose
\begin{align*}
 (u_i-u_j)^\top S_Kv_i
 &= (u_i-u_j)^\top u_i\norm{v_i}^2\\
 &\quad+\mathbf1_{\{j\leq K\}}
 (u_i-u_j)^\top u_j(v_j^\top v_i)\\
 &\quad+\sum_{\substack{k\leq K\\k\notin\{i,j\}}}
 (u_i-u_j)^\top u_k(v_k^\top v_i).
\end{align*}
On \eqref{eq:base-norms}--\eqref{eq:base-inner}, the first term is at
least $4/5$ and the second is $o(1)$.  Conditional only on the anchor
vectors $(u_i,u_j,v_i)$, the summands in the last line are independent,
centered, have total squared conditional $\psi_1$ norm $O(K/d^2)$,
and have maximal conditional $\psi_1$ norm $O(d^{-1})$.  Bernstein's inequality
with $t=(q+4)h$, followed by a union bound over at most $N^3$ triples,
gives
\[
 \left|\sum_{\substack{k\leq K\\k\notin\{i,j\}}}
 (u_i-u_j)^\top u_k(v_k^\top v_i)\right|
 \lesssim \frac{\sqrt{Kh}}d+\frac hd
\]
with failure at most $d^{-q}$.  Choose $\kappa(\beta,q)$ sufficiently
small in \eqref{eq:head-size}; then the last two terms have total
absolute value at most $1/2$ for all large $d$, proving
\eqref{eq:head-margin}.

Choose $C_{\mathrm{mem}}$ so that
$C_{\mathrm{mem}}/4\geq4\alpha+6$.  For $i\leq K$,
\eqref{eq:head-margin} then gives
\[
 \ell_i(W_K)\leq
 \log(1+Ne^{-C_{\mathrm{mem}}h/4})
 \leq N^{-4\alpha-5}.
\]
It remains to control the tail.  Let
\[
 \mathcal F_K:=\sigma(u_1,\ldots,u_N,v_1,\ldots,v_K).
\]
Then $W_K$ is $\mathcal F_K$-measurable, while the vectors
$(v_i)_{i>K}$ are conditionally independent.  For fixed $i>K$, the
function
\[
 f_i(v):=\log\sum_j e^{(u_j-u_i)^\top W_Kv}
\]
is $4\normop{W_K}$-Lipschitz on \eqref{eq:base-norms}.  Conditional
Gaussian concentration therefore gives
\[
 \sum_{i>K}p_i\bigl(f_i(v_i)-\Ee[f_i(v_i)\mid\mathcal F_K]\bigr)
 \lesssim \normop{W_K}\sqrt{\frac{s_Kh}{d}}
\]
simultaneously for all $K$, with failure at most $d^{-q}$.  For the
conditional mean, Jensen's inequality and the Gaussian moment-generating
function give
\[
 \Ee[f_i(v_i)\mid\mathcal F_K]
 \leq h+C\normop{W_K}^2/d.
\]
Multiplying by $p_i$, summing over $i>K$, and using
\eqref{eq:WK-norms} yields \eqref{eq:WK-loss} with one fixed
$C_{\mathrm{cmp}}$.  The head contribution is absorbed by the final
term in \eqref{eq:head-error}.

If \eqref{eq:one-head-window} holds, $W_K$ is feasible by
\eqref{eq:WK-loss}, which gives the operator upper bound.  Its upper
endpoint is exactly the hypothesis of \Cref{lem:pooled}, which gives the
Frobenius lower bound.
\end{proof}

\begin{lemma}[Coverage for $1/2<\alpha<1$]
\label{lem:between-coverage}
Let $1/2<\alpha<1$ and define
\begin{equation}
 K_*:=\left(\frac{d^2}{N^{1-\alpha}h}\right)^{1/\alpha}.
 \label{eq:Kstar-between}
\end{equation}
There is a fixed $\kappa_< >0$ such that, whenever
$K_*/d\to\infty$, the integer
\[
 M:=\lfloor\kappa_< K_*\rfloor
\]
satisfies $M\geq d$ for all large $d$, every $d\leq K\leq M$ obeys
\eqref{eq:head-size} and \eqref{eq:head-admissible}, and the union of
the corresponding gain windows contains
\begin{equation}
 F_d\leq\Delta\leq(1-c_0/4)F_M,
 \qquad
 F_K:=h\sum_{i\leq K}p_i.
 \label{eq:exact-between-covered}
\end{equation}
Consequently,
\begin{equation}
 h(d/N)^{1-\alpha}\lesssim\Delta\lesssim
 h\left(\frac{d^2}{Nh}\right)^{(1-\alpha)/\alpha}
 \label{eq:between-covered-orders}
\end{equation}
is covered.  The condition $K_*/d\to\infty$ is equivalent to
\eqref{eq:between-condition}.
\end{lemma}

\begin{proof}
For $K\leq N/2$, \Cref{lem:zipf} gives
\begin{equation}
 Kp_Kh\asymp N^{\alpha-1}K^{1-\alpha}h,
 \qquad
 s_K\asymp N^{2\alpha-2}K^{1-2\alpha}.
 \label{eq:between-head-scales}
\end{equation}
Also $K_*/N=(d^2/(Nh))^{1/\alpha}\to0$, so all $K\leq M$ are at most
$N/2$ for large $d$.  For $K\geq d$, the deterministic part of
\eqref{eq:head-error} is at most $CKh^2/d^2$.  Dividing by the first
quantity in \eqref{eq:between-head-scales} gives
\[
 C\frac{N^{1-\alpha}K^\alpha h}{d^2}
 \leq C\kappa_<^\alpha.
\]
The fluctuation term divided by $Kp_Kh$ is $O(\sqrt h/d)$, uniformly in
$d\leq K\leq M$, and the final $N^{-4\alpha-4}$ term is uniformly
$o(Kp_Kh)$.  Choose $\kappa_<$ sufficiently small and then $d$
sufficiently large; this proves \eqref{eq:head-admissible}.  Moreover,
\[
 \frac{K_*h}{d^2}
 =\left(\frac{d^2}{Nh}\right)^{(1-\alpha)/\alpha}\to0,
\]
so \eqref{eq:head-size} also holds after decreasing $\kappa_<$.

The gain version of \eqref{eq:one-head-window} contains
\begin{equation}
 I_K:=
 \left[F_K-c_0Kp_Kh,
       F_K-\frac{c_0}{4}Kp_Kh\right].
 \label{eq:between-IK}
\end{equation}
Let $L_K,U_K$ be its two endpoints.  Since $p_{K+1}\leq p_K$,
\[
 (K+1)p_{K+1}-Kp_K\leq p_{K+1}.
\]
Hence
\[
 L_{K+1}-L_K\geq(1-c_0)p_{K+1}h>0,
 \qquad
 U_{K+1}-U_K\geq(1-c_0/4)p_{K+1}h>0.
\]
Thus the intervals move to the right.  Consecutive intervals overlap if
$L_{K+1}\leq U_K$, equivalently if
\begin{equation}
 [c_0(K+1)-1]p_{K+1}
 \geq(c_0/4)Kp_K.
 \label{eq:between-overlap}
\end{equation}
For all sufficiently large $K$,
$c_0(K+1)-1\geq(c_0/2)K$ and
$p_{K+1}/p_K\geq3/4$, so the left side of
\eqref{eq:between-overlap} is at least
$(3c_0/8)Kp_K$.  Therefore
\[
 \bigcup_{K=d}^M I_K=[L_d,U_M].
\]
Because $Kp_Kh\leq F_K$,
\[
 L_d\leq F_d,
 \qquad
 U_M\geq(1-c_0/4)F_M.
\]
Furthermore,
\[
 \frac{F_M}{F_d}\asymp(M/d)^{1-\alpha}\to\infty,
\]
so $F_d\leq(1-c_0/4)F_M$ for all large $d$.  This proves the exact
inclusion \eqref{eq:exact-between-covered}.

Finally,
$F_K\asymp h(K/N)^{1-\alpha}$ and
$M\asymp K_*$, which gives \eqref{eq:between-covered-orders}.  Directly,
\[
 \frac{K_*}{d}
 =d^{[2-\alpha-\beta(1-\alpha)]/\alpha}h^{-1/\alpha},
\]
which tends to infinity exactly under \eqref{eq:between-condition}.
\end{proof}

\begin{lemma}[Coverage at $\alpha=1$]
\label{lem:one-coverage}
Assume $\alpha=1$.  There are fixed constants
$\kappa_1=\kappa_1(\beta,q)>0$ and $A=A(\beta,q)>0$ such that, with
\[
 M:=\left\lfloor\kappa_1\frac{d^2}{h^2}\right\rfloor,
\]
every integer $d\leq K\leq M$ obeys
\eqref{eq:head-size} and \eqref{eq:head-admissible}, and the union of
the corresponding loss windows contains
\begin{equation}
 \log\!\left(\frac{Nh^2}{d^2}\right)+A
 \leq\eps\leq
 \log\!\left(\frac Nd\right)-A.
 \label{eq:one-covered}
\end{equation}
\end{lemma}

\begin{proof}
At $\alpha=1$,
\[
 Kp_Kh=h/H_{N,1}=1+o(1),
 \qquad s_K\asymp(Kh^2)^{-1}.
\]
For $K\geq d$, \eqref{eq:head-error} is bounded by
\[
 CKh^2/d^2+C\sqrt h/d+o(1).
\]
Choose $\kappa_1$ sufficiently small and then $d$ large; this proves
\eqref{eq:head-admissible}.  Also
$Kh\leq\kappa_1d^2/h\leq\kappa d^2$, so
\eqref{eq:head-size} holds.

Write $b_K:=P_Kh$.  Uniformly in this range,
\begin{equation}
 b_K=\log(N/K)+O(1).
 \label{eq:bK-one}
\end{equation}
For all large $d$, $7/8\leq Kp_Kh\leq9/8$, while admissibility gives
$C_{\mathrm{cmp}}\mathsf e_K\leq(c_0/4)Kp_Kh$.  Hence the fixed interval
\[
 J_K:=[b_K+c_0/3,\ b_K+c_0/2]
\]
is contained in the one-head window.  Since
\[
 b_K-b_{K+1}=p_{K+1}h\leq2/(K+1)<c_0/6
\]
for $K\geq d$ and large $d$, consecutive $J_K$ overlap.  Their
endpoints decrease with $K$, so
\[
 \bigcup_{K=d}^M J_K=[b_M+c_0/3,\ b_d+c_0/2].
\]
Using \eqref{eq:bK-one} and
$M\asymp d^2/h^2$ gives
\[
 b_M=\log(Nh^2/d^2)+O(1),
 \qquad b_d=\log(N/d)+O(1).
\]
A sufficiently large fixed $A$ therefore makes
\eqref{eq:one-covered} a subinterval of this union.  Its length tends
to infinity because $\log(d/h^2)\to\infty$.
\end{proof}

\begin{lemma}[Coverage and inversion for $\alpha>1$]
\label{lem:above-coverage}
Assume $\alpha>1$.  There exist a fixed integer $K_0=K_0(\alpha)$, a
constant $\kappa_>=\kappa_>(\alpha,\beta,q)>0$, and
$\eta=\eta(\alpha)\in(0,1)$ such that, with
\[
 M:=\left\lfloor\kappa_>\frac{d^2}{h}\right\rfloor,
\]
every $K_0\leq K\leq M$ obeys \eqref{eq:head-size} and the stronger
bound
\begin{equation}
 C_{\mathrm{cmp}}\mathsf e_K
 \leq\frac{c_0}{16}Kp_Kh.
 \label{eq:above-strong-admissible}
\end{equation}
The corresponding loss windows cover
\begin{equation}
 \frac{h^\alpha}{d^{2(\alpha-1)}}\lesssim\eps\leq\eta h.
 \label{eq:above-covered}
\end{equation}
For every covered $\eps$, one may choose $K$ so that
\begin{equation}
 K\asymp(h/\eps)^{1/(\alpha-1)}.
 \label{eq:above-inversion}
\end{equation}
Finally, the fixed comparator $W_{K_0}$ satisfies
\begin{equation}
 L(W_{K_0})<\eta h.
 \label{eq:fixed-head-comparator}
\end{equation}
\end{lemma}

\begin{proof}
For $K\leq N/2$, \Cref{lem:zipf} gives
\[
 P_K\asymp K^{1-\alpha},
 \qquad Kp_K\asymp K^{1-\alpha},
 \qquad s_K\asymp K^{1-2\alpha}.
\]
Dividing \eqref{eq:head-error} by $Kp_Kh$ yields, uniformly for
$1\leq K\leq M$,
\begin{equation}
 \frac{C_{\mathrm{cmp}}\mathsf e_K}{Kp_Kh}
 \lesssim
 \frac hd+\frac{Kh}{d^2}
 +\sqrt{\frac h{dK}}+\frac{\sqrt h}{d}+o(1).
 \label{eq:above-error-ratio}
\end{equation}
Choose $\kappa_>$ sufficiently small and then $d$ sufficiently large;
this proves \eqref{eq:above-strong-admissible} and
\eqref{eq:head-size}.  Also $M/N\to0$, so all quoted Zipf estimates are
valid.

The one-head window contains
\begin{equation}
 J_K:=h\left[
 P_K+\frac{c_0}{8}Kp_K,
 P_K+\frac{c_0}{2}Kp_K
 \right].
 \label{eq:above-JK}
\end{equation}
Choose $K_0$ so large that, for every $K\geq K_0$,
\begin{equation}
 \left[\frac{c_0}{2}(K+1)-1\right]p_{K+1}
 \geq\frac{c_0}{8}Kp_K.
 \label{eq:above-overlap}
\end{equation}
This is possible because $p_{K+1}/p_K\to1$.  Since
$P_{K+1}=P_K-p_{K+1}$, condition \eqref{eq:above-overlap} is exactly the
statement that the upper endpoint of $J_{K+1}$ is at least the lower
endpoint of $J_K$.  Both endpoints decrease with $K$, because
$P_K$ and $Kp_K=K^{1-\alpha}/H_{N,\alpha}$ decrease.  Thus
\[
 \bigcup_{K=K_0}^M J_K
 =\left[\inf J_M,\ \sup J_{K_0}\right].
\]
Every $\eps\in J_K$ satisfies
\begin{equation}
 \eps\asymp hK^{1-\alpha},
 \label{eq:above-window-scale}
\end{equation}
which proves \eqref{eq:above-inversion}.  Moreover,
\[
 \inf J_M\asymp hM^{1-\alpha}
 \asymp \frac{h^\alpha}{d^{2(\alpha-1)}}.
\]

It remains to choose the fixed upper endpoint.  Let
$p_i^{(\infty)}=i^{-\alpha}/\zeta(\alpha)$ and
$P_K^{(\infty)}=\sum_{i>K}p_i^{(\infty)}$.  For the fixed $K_0$, set
\[
 a:=P_{K_0}^{(\infty)}+\frac{c_0}{16}K_0p_{K_0}^{(\infty)},
 \qquad
 b:=P_{K_0}^{(\infty)}+\frac{c_0}{8}K_0p_{K_0}^{(\infty)},
 \qquad
 \eta:=\frac{a+b}{2}.
\]
Then $0<a<\eta<b<1$: indeed
$Kp_K^{(\infty)}\leq1-P_K^{(\infty)}$ and $c_0/8<1$.
For all large $d$,
\begin{equation}
 P_{K_0}+\frac{c_0}{16}K_0p_{K_0}
 <\eta<
 P_{K_0}+\frac{c_0}{8}K_0p_{K_0}.
 \label{eq:eta-sandwich}
\end{equation}
The right inequality places $\eta h$ below the lower endpoint of
$J_{K_0}$, whereas $\inf J_M=o(h)$.  The connected union therefore
contains \eqref{eq:above-covered} after fixing the implicit lower
constant.  Finally, \eqref{eq:WK-loss},
\eqref{eq:above-strong-admissible}, and the left inequality in
\eqref{eq:eta-sandwich} give \eqref{eq:fixed-head-comparator}.
\end{proof}

\begin{theorem}[Consequences of the head windows]
\label{thm:head-results}
On the intersection of the events in
\Cref{prop:base-event,lem:pooled,prop:head-comparator}, the following
hold simultaneously.
\begin{enumerate}[label=\textup{(\roman*)}]
\item Under the assumptions of \Cref{lem:between-coverage}, every gain
in \eqref{eq:between-covered-orders} satisfies
\[
 \mathcal R_F(\Delta)/\mathcal R_{\mathrm{op}}(\Delta)\asymp\sqrt d.
\]
\item At $\alpha=1$, every loss in \eqref{eq:one-covered} satisfies
$\RF(\eps)/\Rop(\eps)\asymp\sqrt d$.
\item If $\alpha>1$, then every loss in
\eqref{eq:above-covered} satisfies
\begin{equation}
 \frac{\RF(\eps)}{\Rop(\eps)}
 \asymp
 \min\left\{\sqrt d,
 (h/\eps)^{1/(2(\alpha-1))}\right\}.
 \label{eq:above-head-formula}
\end{equation}
The same formula holds for $\eta h\leq\eps<h$.
\end{enumerate}
\end{theorem}

\begin{proof}
For every one-head window, \eqref{eq:one-head-radii} gives
\begin{equation}
 \frac{\RF(\eps)}{\Rop(\eps)}
 \gtrsim\frac{\sqrt K}{1+\sqrt{K/d}}
 \asymp\min\{\sqrt K,\sqrt d\}.
 \label{eq:head-ratio-lower}
\end{equation}
In the first two cases, the covering constructions use $K\geq d$, so
\eqref{eq:head-ratio-lower} and the universal upper bound
\eqref{eq:universal-ratio} give $\asymp\sqrt d$.

Now let $\alpha>1$ and $\eps\leq\eta h$.  The same comparator is
feasible, so \eqref{eq:WK-norms} also gives
$\RF(\eps)\lesssim h\sqrt K$.  Moreover,
$h-\eps\geq(1-\eta)h$ and \eqref{eq:base-nuclear} gives
$\normnuc{G}\asymp1$, so the tangent bound yields
$\Rop(\eps)\gtrsim h$.  Hence
\[
 \frac{\RF(\eps)}{\Rop(\eps)}
 \lesssim\min\{\sqrt K,\sqrt d\}.
\]
Combine this with \eqref{eq:head-ratio-lower} and
\eqref{eq:above-inversion}.

For $\eta h\leq\eps<h$, let $W_0=W_{K_0}$ and
$\Delta_0=h-L(W_0)\gtrsim h$.  Convex scaling gives, with
$\Delta=h-\eps$,
\[
 \RF(\eps)+\Rop(\eps)\lesssim\Delta,
\]
because $K_0$ is fixed and $\normF{W_0}+\normop{W_0}\lesssim h$.
The two tangent bounds and
$\normF{G}\asymp\normnuc{G}\asymp1$ give the reverse inequalities.
Thus the ratio is $\asymp1$, which agrees with
\eqref{eq:above-head-formula} because $h/\eps$ is bounded above and
below by fixed constants.
\end{proof}

% ============================================================
\section{Weak-weight all-accuracy obstruction}
% ============================================================

This section assumes
\begin{equation}
 \frac12<\alpha<1,
 \qquad \beta(1-\alpha)>2.
 \label{eq:weak-regime}
\end{equation}
The companion note~\cite{companion} proves that, because $N>d^2$, a finite minimizer
exists with probability at least
\begin{equation}
 1-2^{1-N}\sum_{r=0}^{d^2-1}\binom{N-1}{r}.
 \label{eq:companion-prob}
\end{equation}

\begin{lemma}[Strict convexity and coercivity]
\label{lem:strict-coercive}
Assume $N\geq d+1$.  Almost surely, $\Psi$ is strictly convex and
coercive on $\R^{d\times d}$.
\end{lemma}

\begin{proof}
Almost surely, $v_1,\ldots,v_d$ span $\R^d$.  Also
\[
 \operatorname{span}\{x_1,\ldots,x_N\}
 =\operatorname{span}\{u_1-u_N,\ldots,u_{N-1}-u_N\}.
\]
The $d\times d$ matrix
$[u_1-u_N,\ldots,u_d-u_N]$ has a nondegenerate Gaussian law on
$\R^{d\times d}$: for coefficient vectors $a_1,\ldots,a_d$ not all
zero,
\[
 \Var\!\left(\sum_{r=1}^d a_r^\top(u_r-u_N)\right)
 =\frac1d\sum_{r=1}^d\norm{a_r}^2
 +\frac1d\norm{\sum_{r=1}^d a_r}^2>0.
\]
Hence its determinant is nonzero almost surely, and the $x_j$ span
$\R^d$.

If $D^2\Psi(W)[A,A]=0$, then by \eqref{eq:hessian-variance}, for each
$i$, the numbers $x_j^\top Av_i$ are constant in $j$.  Their uniform
average is zero, so $x_j^\top Av_i=0$ for all $i,j$.  The two spanning
properties imply $A=0$.  Hence the Hessian is positive definite in
every nonzero direction, proving strict convexity.

We next show that the origin lies in the interior of
$\operatorname{conv}\{x_j\}$.  Otherwise a supporting-hyperplane
argument would give $y\neq0$ such that $x_j^\top y\geq0$ for every
$j$.  Averaging and using $N^{-1}\sum_jx_j=0$ forces
$x_j^\top y=0$ for every $j$, contradicting the spanning property.
Therefore there is a random $r_x>0$ such that
\[
 \max_j x_j^\top y\geq r_x\norm y
 \qquad(y\in\R^d).
\]
For every $y$,
\[
 \log\left(\frac1N\sum_j e^{x_j^\top y}\right)
 \geq\max_jx_j^\top y-h
 \geq r_x\norm y-h.
\]
Applying this with $y=Wv_i$ and averaging over $i$ gives
\[
 \Psi(W)\geq r_x\sum_i p_i\norm{Wv_i}-h.
\]
The map $W\mapsto\sum_i p_i\norm{Wv_i}$ is a norm because the $v_i$
span $\R^d$.  By equivalence of norms on the finite-dimensional matrix
space, it is bounded below by a positive random multiple of
$\normF{W}$.  Hence $\Psi(W)\to\infty$ as $\normF{W}\to\infty$.
\end{proof}

\begin{lemma}[Minimizer homotopy and Fisher sandwich]
\label{lem:homotopy-fisher}
Assume that $F_1(W):=\Psi(W)-\ip{G}{W}$ has a finite minimizer.  For
$t\in[0,1]$, let $W_t$ be the minimizer of
$F_t(W):=\Psi(W)-t\ip{G}{W}$.  Then $W_t$ is uniquely defined,
$t\mapsto W_t$ is $C^1$, $W_0=0$, $W_1=W_\star$, and
\begin{equation}
 \nabla\Psi(W_t)=tG,
 \qquad
 \dot W_t=[\nabla^2\Psi(W_t)]^{-1}[G].
 \label{eq:homotopy}
\end{equation}
Moreover, if
\[
 \rho(W):=\max_{i,j}|x_j^\top Wv_i|\leq\rho_0,
\]
then for every $A$,
\begin{equation}
 e^{-4\rho_0}\tr(A^\top\Sigma_x A M_v)
 \leq D^2\Psi(W)[A,A]
 \leq e^{4\rho_0}\tr(A^\top\Sigma_x A M_v).
 \label{eq:fisher}
\end{equation}
On \eqref{eq:base-Sigma} and \eqref{eq:base-Mv-two-sided}, the last
quantity is $\asymp d^{-2}\normF{A}^2$.
\end{lemma}

\begin{proof}
By \Cref{lem:strict-coercive}, $\Psi$ is strictly convex and coercive.
For $t<1$,
$F_t=(1-t)\Psi+tF_1$ is coercive because $F_1$ is bounded below; at
$t=1$, existence is assumed.  Strict convexity gives uniqueness.  The
first-order condition is $\nabla\Psi(W_t)=tG$.  Since the Hessian is
positive definite, the implicit function theorem gives local $C^1$
branches; uniqueness makes the branches agree on overlaps, yielding the
global path and \eqref{eq:homotopy}.  Since $\Psi\geq0$ and
$\Psi(0)=0$, strict convexity gives $W_0=0$, while $F_1=L-h$ gives
$W_1=W_\star$.

If $\rho(W)\leq\rho_0$, then
$e^{-2\rho_0}/N\leq q_{ji}(W)\leq e^{2\rho_0}/N$.  For any probability
vector $q$,
\[
 \Cov_q(x)=\frac12\sum_{j,k}q_jq_k(x_j-x_k)(x_j-x_k)^\top,
\]
while
\[
 \Sigma_x=\frac1{2N^2}\sum_{j,k}(x_j-x_k)(x_j-x_k)^\top.
\]
Comparison of the coefficients, followed by
\eqref{eq:hessian-variance}, gives \eqref{eq:fisher}.  The final claim
follows from the two-sided covariance bounds.
\end{proof}

\begin{theorem}[Weak-weight obstruction]
\label{thm:weak}
On the intersection of the base Gaussian event and the companion
finite-minimizer event,
\begin{equation}
 \normF{W_\star}\lesssim d^2\sqrt{s_2},
 \qquad
 h-L(W_\star)\gtrsim d^2s_2.
 \label{eq:weak-localization}
\end{equation}
Consequently, simultaneously for every feasible
$L_\star\leq\eps<h$,
\begin{equation}
 \frac{\RF(\eps)}{\Rop(\eps)}\lesssim d^{1-\alpha}.
 \label{eq:weak-obstruction}
\end{equation}
\end{theorem}

\begin{proof}
By \eqref{eq:s2-orders} and \eqref{eq:weak-regime},
\begin{equation}
 d^2\sqrt{s_2}\asymp d^2N^{\alpha-1}\longrightarrow0.
 \label{eq:weak-small}
\end{equation}
Fix a small $\rho_0>0$ and define the exit time
\[
 \tau:=\sup\{t\in[0,1]:\rho(W_s)\leq\rho_0
 \text{ for every }0\leq s\leq t\}.
\]
Because $W_0=0$ and the path is continuous, $\tau>0$.  The defining
property also holds at $\tau$ by continuity.  On $[0,\tau]$,
\eqref{eq:fisher} gives
\[
 \normF{\dot W_t}\lesssim d^2\normF{G},
\]
so
\[
 \normF{W_t}\lesssim d^2\sqrt{s_2}.
\]
On \eqref{eq:base-norms},
$\rho(W_t)\lesssim\normF{W_t}$; by \eqref{eq:weak-small}, this is at
most $\rho_0/2$ for all large $d$.  If $\tau<1$, continuity would
extend the defining property beyond $\tau$, a contradiction.  Hence
$\tau=1$ and the first inequality in \eqref{eq:weak-localization}
holds.

For the gain, take $W=t_0d^2G$ with a sufficiently small fixed $t_0$.
On \eqref{eq:base-norms},
\[
 \max_{0\leq s\leq1}\rho(sW)
 \lesssim\normF{W}\lesssim d^2\sqrt{s_2}=o(1),
\]
so the segment from $0$ to $W$ lies in the same small-logit ball.  The
upper Fisher bound and Taylor's formula give
\[
 L(W)
 \leq h-t_0d^2\normF{G}^2
 +Ct_0^2d^2\normF{G}^2
 \leq h-cd^2s_2.
\]
Optimality of $W_\star$ proves the second inequality in
\eqref{eq:weak-localization}.

Finally, \eqref{eq:all-accuracy}, \eqref{eq:base-nuclear}, and
\eqref{eq:weak-localization} give, for every feasible positive gain,
\[
 \frac{\mathcal R_F(\Delta)}
 {\mathcal R_{\mathrm{op}}(\Delta)}
 \lesssim
 \frac{d^2\sqrt{s_2}\,
 d^{1-\alpha}\sqrt{s_2}}{d^2s_2}
 =d^{1-\alpha}.
\]
\end{proof}

% ============================================================
\section{Assembly of the main theorem}
% ============================================================

\begin{proof}[Proof of \Cref{thm:main}]
Parts (i) and (ii) follow from \Cref{thm:diffuse} and
\eqref{eq:diffuse-mass}.  The local statements in parts (iii) and (iv)
are \Cref{cor:local-phase} together with \eqref{eq:s2-orders}.
The head statement for $1/2<\alpha<1$ follows from
\Cref{lem:between-coverage,thm:head-results}; its nonemptiness was proved
through $K_*/d\to\infty$.  The second statement at $\alpha=1$ follows
from \Cref{lem:one-coverage,thm:head-results}.  Part (v) follows from
\Cref{lem:above-coverage,thm:head-results}; the square-root subinterval
is obtained by comparing
$(h/\eps)^{1/(2(\alpha-1))}$ with $\sqrt d$.  Since its lower endpoint
is $o(h/d^{\alpha-1})$, the asserted additive and multiplicative widths
follow.

The positive statements use a finite intersection of the base event,
the diffuse event when $\alpha\leq1/2$, or the head event when
$\alpha>1/2$.  The corresponding failure probabilities sum to
$Cd^{-q}+Ce^{-cd}$.  All estimates were proved simultaneously over the
required integer indices, and the interval arguments are deterministic,
which proves \eqref{eq:main-prob}.  The final statement is
\Cref{thm:weak}; intersecting with the companion event and using
\eqref{eq:companion-prob} gives \eqref{eq:weak-main-prob}.
\end{proof}

% ============================================================
\section{Subquadratic interpolation regime: \texorpdfstring{$0<\beta<2$}{0 < beta < 2}}
\label{sec:subquadratic}
% ============================================================

In this section only, replace the standing assumption $\beta>2$ by
\begin{equation}
 0<\beta<2,
 \qquad N=\lceil d^\beta\rceil,
 \qquad h=\log N,
 \label{eq:subquadratic-regime}
\end{equation}
while keeping the model, the Zipf weights, and the definitions of
$R_{\mathrm F}(\eps)$ and $R_{\mathrm{op}}(\eps)$ unchanged.  Thus
$Nh=o(d^2)$.  For $x>0$, write
$\log_+x:=\max\{\log x,0\}$.

Set
\begin{equation}
 S:=\sum_{i=1}^N u_iv_i^\top,
 \qquad
 B:=\sum_{i=1}^N x_iv_i^\top,
 \qquad x_i=u_i-\bar u.
 \label{eq:subquadratic-SB}
\end{equation}

\begin{lemma}[Uniform interpolation geometry]
\label{lem:subquadratic-geometry}
Fix $q>0$.  For all sufficiently large $d$, with probability at least
\begin{equation}
 1-Cd^{-q}-Ce^{-cd},
 \label{eq:subquadratic-geometry-prob}
\end{equation}
one has simultaneously
\begin{align}
 &(u_i-u_j)^\top Sv_i\geq\frac14
 &&(1\leq i\neq j\leq N),
 \label{eq:subquadratic-margin}\\
 &\normF{S}\leq C\sqrt N,
 \qquad
 \normop{S}\leq C\left(1+\sqrt{N/d}\right),
 \label{eq:subquadratic-S-norms}\\
 &\normF{B}\leq C\sqrt N,
 \qquad
 \operatorname{rank}(B)\leq\min\{N,d\}.
 \label{eq:subquadratic-B-norms}
\end{align}
The constants may depend on $(\beta,q)$ but not on $d$ or $N$.
\end{lemma}

\begin{proof}
Standard Gaussian norm and inner-product tails, followed by a union
bound over $O(N^2)$ indices, give an event of probability at least
$1-Cd^{-q}-Ce^{-cd}$ on which
\begin{equation}
 \max_i\left|\norm{u_i}^2-1\right|
 +\max_i\left|\norm{v_i}^2-1\right|\leq\frac1{10}
 \label{eq:subquadratic-anchor-norms}
\end{equation}
and
\begin{equation}
 \max_{i\neq j}
 \left(|u_i^\top u_j|+|v_i^\top v_j|\right)
 \leq C\sqrt{h/d}.
 \label{eq:subquadratic-anchor-inner}
\end{equation}
Because $h/d\to0$, the right-hand side of
\eqref{eq:subquadratic-anchor-inner} is at most $1/20$ for all large
$d$ because the constant $C$ is fixed.

Fix $i\neq j$.  Expanding the margin gives
\begin{align}
 (u_i-u_j)^\top Sv_i
 &= (u_i-u_j)^\top u_i\norm{v_i}^2
 \nonumber\\
 &\quad +(u_i-u_j)^\top u_j(v_j^\top v_i)
 \nonumber\\
 &\quad +\sum_{k\notin\{i,j\}}
 (u_i-u_j)^\top u_k(v_k^\top v_i).
 \label{eq:subquadratic-margin-expansion}
\end{align}
On \eqref{eq:subquadratic-anchor-norms}--
\eqref{eq:subquadratic-anchor-inner}, the first term is at least
\[
 (9/10-1/20)(9/10)>3/4,
\]
whereas the absolute value of the second term is at most
\[
 (11/10+1/20)(1/20)<1/16.
\]
Condition on the anchor vectors $(u_i,u_j,v_i)$.  For
$k\notin\{i,j\}$, the random variables
\[
 X_k:=(u_i-u_j)^\top u_k(v_k^\top v_i)
\]
are conditionally independent and centered.  Moreover, the two
Gaussian factors in $X_k$ are conditionally independent and have
conditional subgaussian norms bounded respectively by
$C\norm{u_i-u_j}/\sqrt d$ and $C\norm{v_i}/\sqrt d$.  Hence, on
\eqref{eq:subquadratic-anchor-norms},
\[
 \left\lVert X_k\right\rVert_{\psi_1\mid u_i,u_j,v_i}\leq C/d.
\]
Conditional Bernstein with $t=A h$, where
$A=A(\beta,q)$ is sufficiently large, yields
\begin{equation}
 \left|\sum_{k\notin\{i,j\}}X_k\right|
 \leq C\left(\frac{\sqrt{Nh}}d+\frac hd\right)
 \label{eq:subquadratic-cross-bound}
\end{equation}
except on an event of conditional probability at most $2e^{-Ah}$.
A union bound over fewer than $N^2$ ordered pairs leaves failure at most
$d^{-q}$.  Since $Nh=o(d^2)$, the right-hand side of
\eqref{eq:subquadratic-cross-bound} tends to zero.  Combining the three
terms in \eqref{eq:subquadratic-margin-expansion} proves
\eqref{eq:subquadratic-margin} for all large $d$.

Apply \Cref{lem:bilinear} with all coefficients equal to one.  This
gives \eqref{eq:subquadratic-S-norms} with failure $Ce^{-cd}$.  Finally,
write
\[
 B=S-\bar u\left(\sum_{i=1}^N v_i\right)^\top.
\]
Since $\sqrt N\,\bar u$ and
$N^{-1/2}\sum_i v_i$ are independent copies of
$\cN(0,\Id_d/d)$, Gaussian norm concentration gives
\[
 \norm{\bar u}\lesssim N^{-1/2},
 \qquad
 \norm{\sum_i v_i}\lesssim\sqrt N
\]
with failure $Ce^{-cd}$.  Thus the centering term has Frobenius norm
$O(1)$, and \eqref{eq:subquadratic-B-norms} follows from
\eqref{eq:subquadratic-S-norms}.  The rank bound is deterministic,
because $B=XV^\top$ for the $d\times N$ matrices
$X=[x_1,\ldots,x_N]$ and $V=[v_1,\ldots,v_N]$.
\end{proof}

\begin{theorem}[Sharp full-memory interpolation phase]
\label{thm:subquadratic-interpolation}
Assume \eqref{eq:subquadratic-regime} and fix $q>0$.  There are constants
$c_\star,c,C>0$, with $c_\star$ absolute and $c,C$ allowed to depend on
$(\alpha,\beta,q)$, such that, on an event of probability at least
\begin{equation}
 1-Cd^{-q}-Ce^{-cd},
 \label{eq:subquadratic-main-prob}
\end{equation}
the following statements hold.

First,
\begin{equation}
 L_\star=0,
 \label{eq:subquadratic-Lstar}
\end{equation}
but no finite matrix attains loss zero.  Second, simultaneously for
every
\begin{equation}
 0<\eps\leq c_\star Np_Nh,
 \label{eq:subquadratic-accuracy-phase}
\end{equation}
one has
\begin{align}
 \RF(\eps)
 &\asymp
 \sqrt N\left(h+\log_+\frac{p_N}{\eps}\right),
 \label{eq:subquadratic-RF}\\
 \Rop(\eps)
 &\asymp
 \left(1+\sqrt{N/d}\right)
 \left(h+\log_+\frac{p_N}{\eps}\right).
 \label{eq:subquadratic-Rop}
\end{align}
Consequently,
\begin{equation}
 \frac{\RF(\eps)}{\Rop(\eps)}
 \asymp
 \sqrt{\min\{N,d\}}
 \asymp
 \begin{cases}
 d^{\beta/2},&0<\beta<1,\\
 \sqrt d,&1\leq\beta<2.
 \end{cases}
 \label{eq:subquadratic-separation}
\end{equation}
The upper endpoint of \eqref{eq:subquadratic-accuracy-phase} has the
explicit order
\begin{equation}
 Np_Nh\asymp
 \begin{cases}
 h,&0\leq\alpha<1,\\
 1,&\alpha=1,\\
 hN^{1-\alpha},&\alpha>1.
 \end{cases}
 \label{eq:subquadratic-endpoint}
\end{equation}

every displayed assertion being simultaneous over the whole continuum
of loss levels in \eqref{eq:subquadratic-accuracy-phase}.
\end{theorem}

\begin{proof}
Work on the event in \Cref{lem:subquadratic-geometry}.  By
\eqref{eq:subquadratic-margin}, for every $t\geq0$ and every $i$,
\[
 \ell_i(tS)
 \leq
 \log\left(1+(N-1)e^{-t/4}\right).
\]
The right-hand side tends to zero as $t\to\infty$, and therefore
$L_\star=0$.  On the other hand, if $W$ is finite and $N\geq2$, then
all terms in the defining exponential sum are strictly positive, so
$\ell_i(W)>0$ for every $i$ and hence $L(W)>0$.  Thus the infimum is not
attained at a finite matrix.

Fix a level satisfying \eqref{eq:subquadratic-accuracy-phase}, and put
\begin{equation}
 a_\eps:=h+\log_+\frac{p_N}{\eps}.
 \label{eq:subquadratic-aeps}
\end{equation}
Because $H_{N,\alpha}\leq N$,
\[
 p_N=\frac{N^{-\alpha}}{H_{N,\alpha}}
 \geq N^{-(\alpha+1)},
 \qquad
 \log(1/p_N)\leq(\alpha+1)h.
\]
If $\eps\leq p_N$, then
\[
 \log(N/\eps)
 =h+\log(1/p_N)+\log(p_N/\eps)
 \leq(\alpha+2)a_\eps.
\]
If $\eps>p_N$, then $a_\eps=h$ and again
\[
 \log(N/\eps)
 \leq h+\log(1/p_N)
 \leq(\alpha+2)a_\eps.
\]
Hence, with
\[
 t_\eps:=4(\alpha+2)a_\eps,
\]
one has $(N-1)e^{-t_\eps/4}\leq\eps$, and therefore
$L(t_\eps S)\leq\eps$.  The norm bounds
\eqref{eq:subquadratic-S-norms} give
\begin{equation}
 \RF(\eps)\lesssim\sqrt N\,a_\eps,
 \qquad
 \Rop(\eps)\lesssim
 \left(1+\sqrt{N/d}\right)a_\eps.
 \label{eq:subquadratic-upper-radii}
\end{equation}

We next prove the matching lower bounds.  Let $W$ satisfy
$L(W)\leq\eps$.  Jensen's inequality over all $N$ classes gives, as in
\Cref{lem:pooled},
\begin{equation}
 x_i^\top Wv_i\geq h-\ell_i(W).
 \label{eq:subquadratic-pooled-single}
\end{equation}
Since $p_i\geq p_N$ and all losses are nonnegative,
\[
 p_N\sum_{i=1}^N\ell_i(W)
 \leq\sum_{i=1}^Np_i\ell_i(W)
 =L(W)\leq\eps.
\]
Summing \eqref{eq:subquadratic-pooled-single} therefore yields
\begin{equation}
 \ip{W}{B}\geq Nh-\frac\eps{p_N}.
 \label{eq:subquadratic-pooled-alignment}
\end{equation}
Choose, for definiteness, $c_\star=1/32$.  If
$p_N\leq\eps\leq c_\star Np_Nh$, then $a_\eps=h$ and
\eqref{eq:subquadratic-pooled-alignment} gives
\begin{equation}
 \ip{W}{B}\geq(1-c_\star)Nh\gtrsim Na_\eps.
 \label{eq:subquadratic-alignment-large}
\end{equation}

It remains to treat $0<\eps<p_N$.  Let
\[
 \bar u_{-i}:=\frac1{N-1}\sum_{j\neq i}u_j,
 \qquad
 z_i:=(u_i-\bar u_{-i})^\top Wv_i
 =\frac N{N-1}x_i^\top Wv_i.
\]
Jensen's inequality over the $N-1$ distractors gives
\begin{align*}
 \ell_i(W)
 &=\log\left(1+\sum_{j\neq i}
 e^{(u_j-u_i)^\top Wv_i}\right)\\
 &\geq\log\left(1+(N-1)e^{-z_i}\right).
\end{align*}
Because $p_i\ell_i(W)\leq L(W)\leq\eps$,
$\ell_i(W)\leq\eps/p_i\leq\eps/p_N<1$.  Since
$e^x-1\leq2x$ for $0\leq x\leq1$, inversion of the preceding display
gives
\begin{align}
 z_i
 &\geq\log\left(\frac{N-1}{e^{\eps/p_i}-1}\right)
 \geq\log\left(\frac{(N-1)p_i}{2\eps}\right)
 \nonumber\\
 &\geq\log\left(\frac{(N-1)p_N}{2\eps}\right)
 =a_\eps+\log(1-1/N)-\log2.
 \label{eq:subquadratic-individual-margin}
\end{align}
Since $a_\eps\geq h\to\infty$, the right-hand side is at least
$a_\eps/2$ for all large $d$.  Summing over $i$ and using the identity
for $z_i$ yields
\begin{equation}
 \ip{W}{B}
 =\frac{N-1}{N}\sum_{i=1}^Nz_i
 \geq\frac{N-1}{2}a_\eps
 \gtrsim Na_\eps.
 \label{eq:subquadratic-alignment-small}
\end{equation}
Thus \eqref{eq:subquadratic-alignment-large} and
\eqref{eq:subquadratic-alignment-small} imply, throughout the entire
phase,
\begin{equation}
 \ip{W}{B}\gtrsim Na_\eps.
 \label{eq:subquadratic-master-alignment}
\end{equation}

By \eqref{eq:subquadratic-B-norms},
\[
 \normnuc{B}
 \leq\sqrt{\operatorname{rank}(B)}\normF{B}
 \lesssim\sqrt{N\min\{N,d\}}.
\]
Frobenius duality and nuclear--operator duality applied to
\eqref{eq:subquadratic-master-alignment} now give
\begin{align}
 \normF{W}
 &\gtrsim\sqrt N\,a_\eps,
 \label{eq:subquadratic-lower-F}\\
 \normop{W}
 &\gtrsim
 \sqrt{\frac{N}{\min\{N,d\}}}\,a_\eps
 \asymp
 \left(1+\sqrt{N/d}\right)a_\eps.
 \label{eq:subquadratic-lower-op}
\end{align}
Taking infima over all feasible $W$ and combining with
\eqref{eq:subquadratic-upper-radii} proves
\eqref{eq:subquadratic-RF}--\eqref{eq:subquadratic-Rop}.  Their quotient
is
\[
 \frac{\sqrt N}{1+\sqrt{N/d}}
 \asymp\sqrt{\min\{N,d\}},
\]
which gives \eqref{eq:subquadratic-separation} after substituting
$N=\lceil d^\beta\rceil$.  Finally,
\eqref{eq:subquadratic-endpoint} follows from the Zipf estimates in
\Cref{lem:zipf}.  No union bound over $\eps$ is required, because all
arguments after the Gaussian event are deterministic and hold for every
level simultaneously.
\end{proof}


\begin{thebibliography}{9}
\bibitem{companion}
\emph{Finite Existence and Uniqueness of the Gaussian Softmax
Associative-Memory Minimizer}, companion technical note, 2026.

\bibitem{vershynin}
R.~Vershynin,
\emph{High-Dimensional Probability: An Introduction with Applications
in Data Science}, Cambridge University Press, 2018.

\bibitem{tropp}
J.~A.~Tropp,
\emph{An Introduction to Matrix Concentration Inequalities},
Foundations and Trends in Machine Learning, 2015.

\bibitem{janson}
S.~Janson,
\emph{Gaussian Hilbert Spaces}, Cambridge University Press, 1997.
\end{thebibliography}

\end{document}
